\section*{Capítulo \ref{ch:g8:lenses-images} --- Las lentes y las imágenes}

\begin{solution}{exo:g8:lenses-images:1}
Abultamiento central: convergente --- dobla unos hacia otros los rayos
de un \omterm{def:g6:light-propagation:ray}{haz} paralelo y los recoge en el \omterm{def:g8:lenses-images:focus}{foco}. Hueco central:
divergente --- abre el \omterm{def:g6:light-propagation:ray}{haz} en abanico.
\end{solution}

\begin{solution}{exo:g8:lenses-images:2}
El \omterm{def:g8:lenses-images:focus}{foco} es donde se recoge un \omterm{def:g6:light-propagation:ray}{haz} paralelo después de
la \omterm{def:g8:lenses-images:families}{lente}; la \omterm{def:g8:lenses-images:focus}{distancia focal} $f$ es la distancia de la
\omterm{def:g8:lenses-images:families}{lente} al \omterm{def:g8:lenses-images:focus}{foco}. Una $f$ corta significa mucha curvatura y un doblez
fuerte.
\end{solution}

\begin{solution}{exo:g8:lenses-images:3}
Lanza sobre una pared la \omterm{prop:g5:mirrors-reflection:image}{imagen} de la escena lejana, afínala
moviendo la \omterm{def:g8:lenses-images:families}{lente} y mide de la \omterm{def:g8:lenses-images:families}{lente} a la pared: esa
distancia es $f$. La escena tiene que estar lejos para que sus rayos
lleguen prácticamente paralelos --- y los rayos paralelos, por
definición del \omterm{def:g8:lenses-images:focus}{foco}, se recogen exactamente a $f$.
\end{solution}

\begin{solution}{exo:g8:lenses-images:4}
No concentres nunca la \omterm{def:g1:light-and-shadows:source}{luz} del Sol cerca de nada inflamable --- el
punto focal chamusca y prende en segundos. No mires nunca al Sol a
través de ninguna \omterm{def:g8:lenses-images:families}{lente} --- el ojo más la \omterm{def:g8:lenses-images:families}{lente}
concentran la quemadura sobre la retina, al instante y sin dolor.
\end{solution}

\begin{solution}{exo:g8:lenses-images:5}
Un cuadro montado de verdad por rayos recogidos, que se puede recoger
en una pantalla. Sus señas: invertido y brillante --- la
\omterm{def:g8:lenses-images:families}{lente} cosecha una \omterm{def:g8:lenses-images:families}{lente} entera de \omterm{def:g1:light-and-shadows:source}{luz} por punto allí
donde el agujerito dejaba pasar un hilo.
\end{solution}

\begin{solution}{exo:g8:lenses-images:6}
\omterm{def:g8:lenses-images:families}{Lente}: el objetivo de vidrio --- la \omterm{def:g8:lenses-images:families}{lente} blanda del ojo
detrás de la pupila. Pantalla: el sensor --- la retina. Enfoque: la
cámara desliza su \omterm{def:g8:lenses-images:families}{lente}; el ojo le cambia la forma a la suya, más
gorda para cerca, más delgada para lejos.
\end{solution}

\begin{solution}{exo:g8:lenses-images:7}
La miopía enfoca con demasiada \omterm{def:g2:pushes-and-pulls:force}{fuerza} --- las imágenes se montan
antes de la retina: una \omterm{def:g8:lenses-images:families}{lente divergente} desdobla justo lo
necesario. El ojo hipermétrope que envejece enfoca con demasiada poca
\omterm{def:g2:pushes-and-pulls:force}{fuerza} para el trabajo de cerca: una \omterm{def:g8:lenses-images:families}{lente convergente} presta la
\omterm{def:g2:pushes-and-pulls:force}{fuerza} que falta.
\end{solution}

\begin{solution}{exo:g8:lenses-images:8}
El sello tiene que estar más cerca de la \omterm{def:g8:lenses-images:families}{lente} que su
\omterm{def:g8:lenses-images:focus}{distancia focal}. El ojo disfruta entonces de un fantasma derecho
y agrandado --- rayos doblados como si vinieran de un sello más
grandioso detrás del vidrio; no aterriza nada en ninguna pantalla.
\end{solution}

\begin{solution}{exo:g8:lenses-images:9}
A medida que el objeto se acerca, su \omterm{prop:g5:mirrors-reflection:image}{imagen} retrocede --- así que
la cámara \emph{aleja} su \omterm{def:g8:lenses-images:families}{lente} del sensor para perseguir a la
\omterm{prop:g5:mirrors-reflection:image}{imagen} que se retira. El ojo, en cambio, estira su
\omterm{def:g8:lenses-images:families}{lente} hasta engordarla y dobla más fuerte para traerse la
\omterm{prop:g5:mirrors-reflection:image}{imagen} de vuelta a la retina, que es fija.
\end{solution}

\begin{solution}{exo:g8:lenses-images:10}
A la \omterm{def:g8:lenses-images:focus}{distancia focal}: solo ahí apretuja la \omterm{def:g8:lenses-images:families}{lente} los
rayos prácticamente paralelos del Sol en su punto más pequeño y más
feroz. Más cerca o más lejos, la \omterm{def:g1:light-and-shadows:source}{luz} se reparte por una mancha
demasiado suave para prender --- el incendio espera exactamente en
$f$.
\end{solution}

\begin{solution}{exo:g8:lenses-images:11}
Las gafas de leer son convergentes: recogen la \omterm{def:g1:light-and-shadows:source}{luz} del Sol en un \omterm{def:g8:lenses-images:focus}{foco}
que quema y agrandan la letra de cerca --- dos talentos exclusivos de
las convergentes. Las \omterm{def:g8:lenses-images:families}{lentes divergentes} del nieto miope abren
todos los haces: ni \omterm{def:g8:lenses-images:focus}{foco}, ni punto que queme, y la letra vista a
través de ellas encoge en lugar de crecer.
\end{solution}

\begin{solution}{exo:g8:lenses-images:12}
Telescopio, etapa uno: la gran \omterm{def:g8:lenses-images:families}{lente} recibe los rayos
prácticamente paralelos de una \omterm{def:g4:solar-system:star}{estrella} y los recoge en una
pequeña \omterm{prop:g8:lenses-images:realimage}{imagen real} cerca de su \omterm{def:g8:lenses-images:focus}{foco}; la etapa dos agranda esa
\omterm{prop:g5:mirrors-reflection:image}{imagen}. La \omterm{def:g1:light-and-shadows:source}{luz} débil de las \omterm{def:g4:solar-system:star}{estrellas} es preciosa --- la
\emph{anchura} de la \omterm{def:g8:lenses-images:families}{lente} delantera decide cuánta se cosecha por
\omterm{def:g4:solar-system:star}{estrella}: la abertura lo es todo. Microscopio, etapa uno: una
\omterm{def:g8:lenses-images:families}{lente} muy curvada y de $f$ corta, colgada muy cerca de la gotita,
lanza una \omterm{prop:g8:lenses-images:realimage}{imagen real} enormemente agrandada; la etapa dos vuelve a
agrandar. Su objeto es brillante y cercano --- lo que la etapa uno
necesita no es anchura, sino poder de doblar. La misma arquitectura,
apetitos opuestos.
\end{solution}

\begin{solution}{pb:g8:lenses-images:1}
\textbf{1.} Tacto: centro abultado, convergente; centro hueco,
divergente. Prueba de la ventana: las \omterm{def:g8:lenses-images:families}{lentes convergentes} lanzan
sobre la pared una ventana del revés a cierta distancia; las
\omterm{def:g8:lenses-images:families}{lentes divergentes} no lo hacen nunca.
\textbf{2.} $f = \qty{25}{cm}$; convergente.
\textbf{3.} Divergente --- ninguna \omterm{prop:g8:lenses-images:realimage}{imagen real} a ninguna distancia,
y la vista encogida a través de ella sella el veredicto.
\textbf{4.} La C: la \omterm{def:g8:lenses-images:focus}{distancia focal} más corta (\qty{10}{cm})
señala a la \omterm{def:g8:lenses-images:families}{lente} más gorda y que dobla con más \omterm{def:g2:pushes-and-pulls:force}{fuerza}.
\textbf{5.} Borroso de lejos y cómodo de cerca: miopía --- el ojo dobla
con demasiada \omterm{def:g2:pushes-and-pulls:force}{fuerza}. Receta la familia divergente.
\textbf{6.} Lectura con el brazo estirado: el ojo hipermétrope que
envejece y enfoca con poca \omterm{def:g2:pushes-and-pulls:force}{fuerza}. Receta cristales convergentes de
leer.
\textbf{7.} Los cristales de leer añaden un poder de doblar
presupuestado para las páginas cercanas; apuntados a las montañas,
cuyos rayos llegan paralelos y no necesitan ayuda, esa \omterm{def:g2:pushes-and-pulls:force}{fuerza} añadida
dobla de más --- las imágenes se montan antes de la retina y los
picos se emborronan.
\textbf{8.} Añadirle corrección a un ojo bien enfocado lo desequilibra:
la \omterm{def:g8:lenses-images:families}{lente} de prueba dobla una \omterm{def:g1:light-and-shadows:source}{luz} contra la que el ojo de la niña
tiene entonces que pelear, y monta las imágenes fuera de la
retina --- los paneles que se van son la protesta de un ojo sano.
\textbf{9.} La \omterm{def:g8:lenses-images:families}{lente} delantera recoge los rayos paralelos de la
\omterm{def:g4:solar-system:star}{estrella} en una pequeña \omterm{prop:g8:lenses-images:realimage}{imagen real} cerca de su \omterm{def:g8:lenses-images:focus}{foco}; la
\omterm{def:g8:lenses-images:families}{lente} ocular es una lupa colgada a menos de su propia
\omterm{def:g8:lenses-images:focus}{distancia focal} de esa \omterm{prop:g5:mirrors-reflection:image}{imagen}, que la agranda hacia el
ojo del observador.
\textbf{10.} Cada \omterm{def:g4:solar-system:star}{estrella} entrega apenas un susurro de
\omterm{def:g1:light-and-shadows:source}{luz} por \omterm{def:g2:measuring-length:units}{centímetro} cuadrado; la \omterm{def:g8:lenses-images:families}{lente} ancha
como un platillo recoge \omterm{def:g2:measuring-length:units}{centímetros} a centenares y encauza todo
un plato de susurro hacia un único punto brillante. La anchura es el apetito del
astrónomo.
\textbf{11.} La lupa del tamaño de una moneda, por muy curvada que
esté, cosecharía \omterm{def:g1:light-and-shadows:source}{luz} de \omterm{def:g4:solar-system:star}{estrella} por valor de una moneda allí donde
hace falta un platillo: imágenes quizá nítidas, pero muertas de
hambre --- adiós a las \omterm{def:g4:solar-system:star}{estrellas} débiles. Se gana: nada que no
diera ya la \omterm{def:g8:lenses-images:families}{lente} ocular. Veredicto: no hay recambio; encarga un
vidrio como es debido.
\textbf{12.} Por ejemplo: ``Todo lo que hay en este estante hace dos
verbos: \emph{recoger} \omterm{def:g1:light-and-shadows:source}{luz} y después \emph{entregarla} donde el
ojo la necesita --- y todo el instrumental óptico de la civilización, de
las gafas a los observatorios, son esos dos verbos en tamaños
distintos''.
\end{solution}
